% GENERATED FILE. Edit canonical lessons, then run npm run generate:books.
% canonical-source-hash: fnv1a64:7805c6f437ad8c98
% canonical-lessons: ZH-C491-jiaxin, ZH-C491-shangsi, ZH-W491-si-charge, ZH-C491-mishu, ZH-C491-zhichang, ZH-C491-xiashu, ZH-R491-review-486-491

\chapter{The Boss}
\label{ch:zh-the-boss}
\hlchaptermodality{491}

\begin{chapteropening}
\textbf{By the end of this chapter:} \emph{I can talk about a pay rise, my superior, a secretary, the world of work and the people under me.}

Complete the last lesson of chapter 491: I can talk about a pay rise, my superior, a secretary, the world of work and the people under me.
\end{chapteropening}

\section[jiāxīn]{jiāxīn — to get a pay rise; to raise someone's pay}

\label{lesson:ZH-C491-jiaxin}



\begin{quote}
Before the new one: say the Mandarin for exhausting (of work), then the Mandarin for to fire (from a job).
\end{quote}

\subsection*{You'll want to know: jiāxīn}
\emph{jiāxīn} — \textquotedblleft{}to get a pay rise; to raise someone's pay\textquotedblright{}.

\begin{quote}
\emph{wǒ gōngzuò le liǎng nián hái méiyǒu jiāxīn} — I've worked for two years and still haven't had a pay rise.
\end{quote}

Its characters come later; for now it is a word for the ear.

\begin{cousinweb}
\emph{jiāxīn} is built from \emph{jiā} 'to add' and \emph{xīn} 'firewood', both inherited from Middle Chinese, \emph{kae} and \emph{sin} (Guangyun, in Baxter's transcription), which glosses the first 'to increase' and the second 'firewood'. Firewood standing for pay is the usual explanation, but the dictionary does not give that sense.
\end{cousinweb}

\begin{grammarlens}[title={What you've built}]
One more word for saying how it was.
\end{grammarlens}

\subsection*{Your turn}
\begin{itemize}
\raggedright
  \item \emph{Say it:} \emph{jiāxīn}
  \item \emph{Say it:} \emph{jiāxīn}, once more
  \item \emph{Say it:} \emph{wǒ gōngzuò le liǎng nián hái méiyǒu jiāxīn} — I've worked for two years and still haven't had a pay rise.
  \item \emph{Recall:} say the Mandarin for exhausting (of work), then the Mandarin for to fire (from a job), then say \emph{jiāxīn} again
\end{itemize}

\begin{tcolorbox}[breakable,colback=teal!4,colframe=teal!35!black,title={Before you move on}]
What does jiāxīn mean? (\textquotedblleft{}To get a pay rise; to raise someone's pay\textquotedblright{}.) Say it once more.
\end{tcolorbox}
\section[\texorpdfstring{\zh{上司} (shàngsi)}{上司 (shàngsi)}]{\zh{上司} (shàngsi) — one's superior at work}

\label{lesson:ZH-C491-shangsi}



\begin{quote}
Before the new one: say the Mandarin for to fire (from a job), then the Mandarin for to get a pay rise.
\end{quote}

\subsection*{You'll want to know: \zh{上司}}
\textbf{\zh{上司}} — \emph{shàngsi} — \textquotedblleft{}one's superior at work\textquotedblright{}.

\begin{quote}
\textbf{\zh{我的新上司比以前的好}} \emph{wǒ de xīn shàngsi bǐ yǐqián de hǎo} — My new boss is better than the last one.
\end{quote}

The new character, in pieces: \zh{司} — five strokes: one long stroke for the frame, a bar along the top that turns down the right side and hooks; a short bar under it; then \zh{口} \emph{kǒu} inside. The next lesson puts it in your hand.

The second syllable is light: shàngsi, with no tone of its own.

\begin{cousinweb}
\textbf{\zh{上司}} (\emph{shàngsi}) is built from \zh{上} 'above' and \zh{司} 'to be in charge', both inherited from Middle Chinese, \emph{dzyangH} and \emph{si} (Guangyun, in Baxter's transcription), which glosses the first 'a ruler' and the second 'to be in charge'.
\end{cousinweb}

\begin{grammarlens}[title={What you've built}]
One more word for saying how it was.
\end{grammarlens}

\subsection*{Your turn}
\begin{itemize}
\raggedright
  \item \emph{Say it:} \emph{shàngsi}
  \item \emph{Say it:} \emph{shàngsi}, once more
  \item \emph{Say it:} \emph{wǒ de xīn shàngsi bǐ yǐqián de hǎo} — My new boss is better than the last one.
  \item \emph{Recall:} say the Mandarin for to fire (from a job), then the Mandarin for to get a pay rise, then say \emph{shàngsi} again
\end{itemize}

\begin{tcolorbox}[breakable,colback=teal!4,colframe=teal!35!black,title={Before you move on}]
What does \zh{上司} mean? (\textquotedblleft{}One's superior at work\textquotedblright{}.) Say it once more.
\end{tcolorbox}
\section[\texorpdfstring{\zh{司} (sī)}{司 (sī)}]{\zh{司} — a hooked frame, a short bar, then kǒu}

\label{lesson:ZH-W491-si-charge}



\begin{quote}
The word first, then the character inside it.

\begin{itemize}
\raggedright
  \item \emph{Recall:} say \emph{shàngsi} — \textbf{R1}, one lesson back, before it decays
\end{itemize}

\begin{quote}
\zh{口} \zh{公}
\end{quote}

\emph{Read it:} each one aloud, with nothing said first — \emph{kǒu}, \emph{gōng}
\end{quote}

\subsection*{Script — a hooked frame, a short bar, then kǒu}
\hlblockfigure{\detokenize{figures/ZH-W491-si-charge-filmstrip.pdf}}{How \zh{司} is written, stroke by stroke}

\begin{quote}
\zh{司}
\end{quote}

\textbf{The frame, one stroke.} A bar along the top that turns down the right side and hooks to the left at the foot.

\textbf{A short bar} inside, under the top.

\textbf{Then \zh{口} \emph{kǒu}, three strokes}, under the short bar: the left side, the top that turns down, and the bottom.

1 + 1 + 3 = five.

\subsection*{Your turn}
\begin{itemize}
\raggedright
  \item \emph{Write it:} \zh{司}, all five, with the model in front of you
  \item \emph{Cover:} write it again from memory, and count the pieces, not the strokes
\end{itemize}

\begin{tcolorbox}[breakable,colback=teal!4,colframe=teal!35!black,title={Before you move on}]
How many strokes does \zh{司} have? (\textbf{Five}.) Which stroke comes first? (\textbf{The frame: a top bar that turns down and hooks}.) What sits inside? (\textbf{A short bar and \zh{口}}.)

Source: \href{https://raw.githubusercontent.com/chanind/hanzi-writer-data/68d10a4b21150cae5e1ebbd223eed289cf32d90c/data/\%E5\%8F\%B8.json}{Hanzi Writer Data \zh{司}, pinned stroke paths and medians}.
\end{tcolorbox}
\section[mìshū]{mìshū — secretary}

\label{lesson:ZH-C491-mishu}



\begin{quote}
Before the new one: say the Mandarin for to get a pay rise, then the Mandarin for one's superior at work.
\end{quote}

\subsection*{You'll want to know: mìshū}
\emph{mìshū} — \textquotedblleft{}secretary\textquotedblright{}.

\begin{quote}
\emph{wǒ zài gōngsī dāng le sān nián mìshū} — I was a secretary at a company for three years.
\end{quote}

Its characters come later; for now it is a word for the ear.

\begin{cousinweb}
\emph{mìshū} is built from \emph{mì} 'secret' and \zh{书} 'writing', both inherited from Middle Chinese, \emph{pijH} and \emph{syo} (Guangyun, in Baxter's transcription), which lists the first in an older form, glossed 'close' and 'secret'. The modern mì is irregular: \emph{pijH} would regularly give bì.
\end{cousinweb}

\begin{grammarlens}[title={What you've built}]
One more word for saying how it was.
\end{grammarlens}

\subsection*{Your turn}
\begin{itemize}
\raggedright
  \item \emph{Say it:} \emph{mìshū}
  \item \emph{Say it:} \emph{mìshū}, once more
  \item \emph{Say it:} \emph{wǒ zài gōngsī dāng le sān nián mìshū} — I was a secretary at a company for three years.
  \item \emph{Recall:} say the Mandarin for to get a pay rise, then the Mandarin for one's superior at work, then say \emph{mìshū} again
\end{itemize}

\begin{tcolorbox}[breakable,colback=teal!4,colframe=teal!35!black,title={Before you move on}]
What does mìshū mean? (\textquotedblleft{}Secretary\textquotedblright{}.) Say it once more.
\end{tcolorbox}
\section[\texorpdfstring{\zh{职场} (zhíchǎng)}{职场 (zhíchǎng)}]{\zh{职场} (zhíchǎng) — the workplace; the world of work}

\label{lesson:ZH-C491-zhichang}



\begin{quote}
Before the new one: say the Mandarin for one's superior at work, then the Mandarin for a secretary.
\end{quote}

\subsection*{You'll want to know: \zh{职场}}
\textbf{\zh{职场}} — \emph{zhíchǎng} — \textquotedblleft{}the workplace; the world of work\textquotedblright{}.

\begin{quote}
\emph{wǒ juéde zhíchǎng bǐ xuéxiào lèi} — I find working life more tiring than school.
\end{quote}

\begin{cousinweb}
\textbf{\zh{职场}} (\emph{zhíchǎng}) is built from \zh{职} 'a post' and \zh{场} 'a ground', both inherited from Middle Chinese, \emph{tsyik} and \emph{drjang} (Guangyun, in Baxter's transcription), which glosses the second 'a threshing floor'; the modern third tone of chǎng is irregular for \emph{drjang}.
\end{cousinweb}

\begin{grammarlens}[title={What you've built}]
One more word for saying how it was.
\end{grammarlens}

\subsection*{Your turn}
\begin{itemize}
\raggedright
  \item \emph{Say it:} \emph{zhíchǎng}
  \item \emph{Say it:} \emph{zhíchǎng}, once more
  \item \emph{Say it:} \emph{wǒ juéde zhíchǎng bǐ xuéxiào lèi} — I find working life more tiring than school.
  \item \emph{Recall:} say the Mandarin for one's superior at work, then the Mandarin for a secretary, then say \emph{zhíchǎng} again
\end{itemize}

\begin{tcolorbox}[breakable,colback=teal!4,colframe=teal!35!black,title={Before you move on}]
What does \zh{职场} mean? (\textquotedblleft{}The workplace; the world of work\textquotedblright{}.) Say it once more.
\end{tcolorbox}
\section[xiàshǔ]{xiàshǔ — subordinate; someone under one at work}

\label{lesson:ZH-C491-xiashu}



\begin{quote}
Before the new one: say the Mandarin for a secretary, then the Mandarin for the world of work.
\end{quote}

\subsection*{You'll want to know: xiàshǔ}
\emph{xiàshǔ} — \textquotedblleft{}subordinate; someone under one at work\textquotedblright{}.

\begin{quote}
\emph{tā de xiàshǔ dōu hěn yǒu jīngshen} — The people under him are all full of energy.
\end{quote}

Its characters come later; for now it is a word for the ear.

\begin{cousinweb}
\emph{xiàshǔ} is built from \zh{下} 'below' and \emph{shǔ} 'to belong to', both inherited from Middle Chinese, \emph{haeX} and \emph{dzyowk} (Guangyun, in Baxter's transcription), which glosses the first 'low' and the second 'attached to' and 'a kind'.
\end{cousinweb}

\begin{grammarlens}[title={What you've built}]
One more word for saying how it was.
\end{grammarlens}

\subsection*{Your turn}
\begin{itemize}
\raggedright
  \item \emph{Say it:} \emph{xiàshǔ}
  \item \emph{Say it:} \emph{xiàshǔ}, once more
  \item \emph{Say it:} \emph{tā de xiàshǔ dōu hěn yǒu jīngshen} — The people under him are all full of energy.
  \item \emph{Recall:} say the Mandarin for a secretary, then the Mandarin for the world of work, then say \emph{xiàshǔ} again
\end{itemize}

\begin{tcolorbox}[breakable,colback=teal!4,colframe=teal!35!black,title={Before you move on}]
What does xiàshǔ mean? (\textquotedblleft{}Subordinate; someone under one at work\textquotedblright{}.) Say it once more.
\end{tcolorbox}
\section[(dialogue)]{Looking back: from Looking for Work to The Boss}

\label{lesson:ZH-R491-review-486-491}



\begin{quote}
Say the Mandarin for to recruit staff and a CV.
\end{quote}

\subsection*{Your turn}
\begin{itemize}
\raggedright
  \item \emph{Say it:} to recruit staff, a CV, to apply for a job, to do an internship, pay and conditions
  \item \emph{Say it:} a line of business, one's life's work, to find employment, to be out of work, a factory
  \item \emph{Say it:} an engineer, a chef, a scientist, to translate, to earn money
  \item \emph{Say it:} to move to another employer, a member of staff, full-time (of a job), a position in a firm, to get promoted
  \item \emph{Say it:} to manage, to work together, the night shift, exhausting (of work), to fire (from a job)
  \item \emph{Say it:} to get a pay rise, one's superior at work, a secretary, the world of work, someone under one at work
  \item \emph{Read it:} \zh{业} \zh{职} \zh{司} — each one aloud, then a word that holds it
\end{itemize}

\begin{tcolorbox}[breakable,colback=teal!4,colframe=teal!35!black,title={Before you move on}]
To recruit staff? (\emph{zhāopìn}.) A CV? (\emph{jiǎnlì}.) To apply for a job? (\emph{yìngpìn}.) To do an internship? (\emph{shíxí}.) Pay and conditions? (\emph{dàiyù}.) A line of business? (\emph{hángyè}.) One's life's work? (\textbf{\zh{事业}}, \emph{shìyè}.) To find employment? (\textbf{\zh{就业}}, \emph{jiùyè}.) To be out of work? (\emph{shīyè}.) A factory? (\emph{gōngchǎng}.) An engineer? (\emph{gōngchéngshī}.) A chef? (\emph{chúshī}.) A scientist? (\emph{kēxuéjiā}.) To translate? (\emph{fānyì}.) To earn money? (\emph{zhuàn}.) To move to another employer? (\emph{tiàocáo}.) A member of staff? (\textbf{\zh{职工}}, \emph{zhígōng}.) Full-time (of a job)? (\textbf{\zh{全职}}, \emph{quánzhí}.) A position in a firm? (\emph{zhíwèi}.) To get promoted? (\emph{shēngzhí}.) To manage? (\emph{guǎnlǐ}.) To work together? (\emph{hézuò}.) The night shift? (\emph{yèbān}.) Exhausting (of work)? (\emph{xīnkǔ}.) To fire (from a job)? (\emph{kāichú}.) To get a pay rise? (\emph{jiāxīn}.) One's superior at work? (\textbf{\zh{上司}}, \emph{shàngsi}.) A secretary? (\emph{mìshū}.) The world of work? (\textbf{\zh{职场}}, \emph{zhíchǎng}.) Someone under one at work? (\emph{xiàshǔ}.)
\end{tcolorbox}
